\FloatBarrier
\newpage
\section{随机变量}
\label{sec:prob-random-variables}

同一次试验可以记录不同的数值。抛两次硬币，可以记录正面总数，也可以记录第一次是否正面。数值由实际结果确定，于是可以使用函数的语言来研究它。

\subsection{随机变量的定义}

\begin{definition}
在有限样本空间\(\Omega\)上，一个从样本点到实数的函数\(X:\Omega\to\mathbb R\)称为{\heiti 随机变量}。
\end{definition}

随机的是试验结果，函数\(X\)则把每个结果对应到一个确定的数。例如在两次抛币中定义正面次数\(Y\)，便有
\[
 Y(HH)=2,\quad Y(HT)=Y(TH)=1,\quad Y(TT)=0.
\]
\(\{Y=1\}\)表示事件\(\{HT,TH\}\)，\(\{Y\ge1\}\)表示\(\{HH,HT,TH\}\)。一般地，\(\{X\in A\}\)是所有满足\(X(\omega)\in A\)的样本点组成的事件，它的概率由原模型确定。

若对数值再作运算，\(g(X)\)仍是样本点的函数，也就是新的随机变量。两个变量的和\(X+Y\)、乘积\(XY\)，同样按每个样本点上的函数值定义。

\subsection{分布列}

\begin{definition}
有限随机变量\(X\)的不同可能取值为\(x_1,\ldots,x_m\)，记
\[
 p_j=P(X=x_j),\qquad j=1,\ldots,m.
\]
取值与相应概率组成的表称为\(X\)的{\heiti 分布列}。
\end{definition}

事件\(\{X=x_j\}\)分割样本空间，故\(p_j\ge0\)、\(\sum_jp_j=1\)。两次独立公平抛币的正面次数满足
\[
\begin{array}{c|ccc}
Y&0&1&2\\\hline
P&1/4&1/2&1/4
\end{array}
\]
分布列保留了各数值的概率，却不再逐一记录这些数值来自哪些样本点。不同试验中的变量可以有相同分布；同一空间中的两个变量也可以有相同分布而取值并不处处相同。

求\(g(X)\)的分布时，把产生同一新数值的概率合并。例如\(Z=(Y-1)^2\)只取\(0,1\)，两者的概率都为\(1/2\)。

\begin{example}
公平掷两枚骰子，\(X\)为较大点数。求\(X\)与\(6-X\)的分布列。
\end{example}

\newpage
\subsection{分布函数}

\begin{definition}
随机变量\(X\)的{\heiti 分布函数}定义为
\begin{equation}\label{eq:prob-cdf}
 F_X(x)=P(X\le x),\qquad x\in\mathbb R.
\end{equation}
\end{definition}

它记录从左侧累积到\(x\)的概率，因而单调不减，取值在\([0,1]\)中。对有限变量，在最小可能值左侧为0，在最大可能值及其右侧为1。例如上述\(Y\)满足
\[
 F_Y(x)=\begin{cases}
 0,&x<0,\\1/4,&0\le x<1,\\3/4,&1\le x<2,\\1,&x\ge2.
 \end{cases}
\]

\begin{figure}[!ht]
\centering
\includegraphics[width=.55\linewidth]{probability-cdf}
\caption{两次独立公平抛币的正面次数的分布函数。实心端点表示该处实际函数值。}
\end{figure}

在相邻取值之间，事件\(\{X\le x\}\)不变，图像便保持水平。由于定义包含端点，有限变量的分布函数在跳跃处取跳跃后的高度，并且右连续；图中的实心与空心端点正反映了这一点。

由事件分解，对\(a<b\)总有
\begin{equation}\label{eq:prob-cdf-interval}
 P(a<X\le b)=F_X(b)-F_X(a).
\end{equation}
对离散变量，跳跃大小给出该点的概率：\(P(X=a)=F_X(a)-F_X(a^-)\)，其中\(F_X(a^-)\)为左极限。有限模型中的这一累计表达，后面还可以用于连续取值的模型。

\begin{example}
利用上述\(F_Y\)，求\(P(Y<1)\)、\(P(Y\le1)\)与\(P(1\le Y\le2)\)，再从各跳跃高度恢复\(Y\)的分布列。
\end{example}

\subsection{联合分布}

研究两个变量的共同变化，需要知道它们同时取值的概率。若\(X\)取\(x_i\)、\(Y\)取\(y_j\)，记
\[
 p_{ij}=P(X=x_i,Y=y_j).
\]
这些概率组成{\heiti 联合分布列}，各项非负，总和为1。对同一行或列求和，得到各变量自身的分布，称为{\heiti 边际分布}。

在两次独立公平抛币中，令\(X\)记录第一次是否正面，\(Y\)记录正面总次数，便有
\[
\begin{array}{c|ccc|c}
 X\backslash Y&0&1&2&P(X=x)\\\hline
 0&1/4&1/4&0&1/2\\
 1&0&1/4&1/4&1/2\\\hline
 P(Y=y)&1/4&1/2&1/4&1
\end{array}
\]
已知\(Y=1\)，将这一列的概率除以列和\(1/2\)，就得到\(X\)在该条件下的分布。已知\(Y=2\)时，\(X=1\)的条件概率则为1。

由联合表也可以求\(X+Y\)等变量的分布：先标记每个格子的函数值，再把相同值对应的概率相加。这仍是样本结果合并的同一过程。

\subsection{随机变量的独立性}

\begin{definition}
若对全部可能取值均有
\begin{equation}\label{eq:prob-variable-independence}
 P(X=x_i,Y=y_j)=P(X=x_i)P(Y=y_j),
\end{equation}
称有限随机变量\(X,Y\)相互独立。
\end{definition}

对任意取值集合\(A,B\)，把上式在相应格子中求和，得到
\[
 P(X\in A,Y\in B)=P(X\in A)P(Y\in B).
\]
所以这个定义与事件独立性一致，也说明分别作用于\(X,Y\)的函数\(g(X),h(Y)\)仍然独立。

多个随机变量相互独立，要求任取有限子组及各自的取值，相应联合概率均等于各边际概率的乘积。连续取值时，则用各变量分别落入任意区间的事件满足同样的乘积分解来定义独立性。

两次独立抛币各自的正面指示量独立；第一次的正面指示量与总正面次数却不独立。即使给定两张相同的边际分布，也不能由此确定联合分布。例如\(X,Y\)均以\(1/2\)的概率取0和1，若\(Y=X\)，和只取0、2；若二者独立，和还以\(1/2\)的概率取1。

\begin{exercise}
\mbox{}\par\nobreak

（1）随机变量\(X\)在\(-1,0,1\)处的概率分别为\(1/4,1/2,1/4\)。写出\(X^2\)和\(2X+1\)的分布列及\(X\)的分布函数。

（2）从编号1至4的球中不放回抽两个，令\(X\)为第一次的数字，\(Y\)为第二次的数字。写出联合分布，求\(P(X+Y=5)\)，并判断独立性。

（3）在本节抛币联合表中，求\(X+Y\)的分布及\(P(X=1\mid Y\ge1)\)。

（4）设\(X,Y\)独立，各以\(1/2\)的概率取\(-1,1\)。求\(X+Y\)、\(XY\)的分布，判断\(X\)与\(XY\)是否独立。
\end{exercise}
